\documentclass[12pt,a4paper]{article}
\usepackage[utf8]{inputenc}
\usepackage[brazil]{babel}
\usepackage{amsmath,amssymb,amsfonts}
\usepackage{graphicx}
\usepackage{geometry}
\usepackage{booktabs}
\usepackage{algorithm}
\usepackage{algpseudocode}
\usepackage{hyperref}
\usepackage{xcolor}
\usepackage{listings}
\usepackage{tikz}
\usetikzlibrary{shapes,arrows,positioning}

\geometry{margin=2.5cm}

\title{\textbf{Algoritmos de Alocação: Sequencial e Otimizado}\\
\large Sistema SIVIRA - Otimização de Produção}
\author{Documentação Técnica}
\date{Dezembro 2025}

\begin{document}

\maketitle

\tableofcontents
\newpage

\section{Introdução}

Este documento descreve os dois algoritmos de alocação de pedidos disponíveis no sistema SIVIRA:

\begin{enumerate}
    \item \textbf{Algoritmo Sequencial}: Execução com heurística de priorização
    \item \textbf{Otimizador PL}: Otimização unificada via MILP que suporta ambos os modos:
    \begin{itemize}
        \item Modo DETERMINÍSTICO ($\tau_{max} = 0$): Atividades contíguas
        \item Modo FLEXÍVEL ($\tau_{max} > 0$): Janelas temporais flexíveis
    \end{itemize}
\end{enumerate}

O Otimizador PL utiliza o solver \textbf{OR-Tools CP-SAT} do Google, que oferece alta performance para problemas de satisfação de restrições e otimização combinatória.

\subsection{Objetivos do Sistema}

\begin{itemize}
    \item Maximizar o número de pedidos atendidos dentro dos prazos (deadlines)
    \item Respeitar restrições de equipamentos (capacidade finita)
    \item Garantir sequenciamento correto de atividades (gaps temporais)
    \item Otimizar a ordem de execução considerando urgência dos pedidos
\end{itemize}

\section{Algoritmo Sequencial}

O Algoritmo Sequencial é a abordagem mais simples para execução de pedidos, utilizando uma \textbf{heurística de priorização} baseada em critérios de restrições.

\subsection{Descrição}

O método sequencial ordena os pedidos por uma heurística de priorização antes da execução, considerando: (1) duração total estimada (maiores primeiro), (2) quantidade (maiores primeiro), (3) tempo de antecipação (menores primeiro) e (4) ID do pedido como desempate. É uma abordagem \textbf{gulosa (greedy)} que processa cada pedido individualmente, sem otimização global de recursos.

\subsection{Arquitetura do Algoritmo Sequencial}

\begin{center}
\begin{tikzpicture}[node distance=1.2cm, auto,
    fase/.style={draw, rectangle, rounded corners, minimum width=5cm, minimum height=0.9cm, font=\small}]

    \node[fase, fill=blue!20] (f1) {Fase 1: Receber lista de pedidos};
    \node[fase, fill=green!20, below=of f1] (f2) {Fase 2: Gerar comanda do pedido};
    \node[fase, fill=orange!20, below=of f2] (f3) {Fase 3: Criar atividades modulares};
    \node[fase, fill=red!20, below=of f3] (f4) {Fase 4: Executar atividades em ordem};

    \draw[->, thick] (f1) -- (f2);
    \draw[->, thick] (f2) -- (f3);
    \draw[->, thick] (f3) -- (f4);
    \draw[->, thick, dashed] (f4.east) to[out=0,in=0] node[right] {próximo pedido} (f2.east);
\end{tikzpicture}
\end{center}

\subsection{Pseudocódigo}

\begin{algorithm}
\caption{Execução Sequencial de Pedidos}
\begin{algorithmic}[1]
\State $pedidos\_executados \gets 0$
\State $pedidos\_com\_erro \gets 0$
\For{cada $pedido$ em $lista\_pedidos$}
    \State Gerar comanda com reserva de ingredientes
    \State Criar atividades modulares necessárias
    \State Executar atividades em ordem cronológica
    \If{execução bem-sucedida}
        \State $pedidos\_executados \gets pedidos\_executados + 1$
        \State Salvar log de sucesso
    \Else
        \State $pedidos\_com\_erro \gets pedidos\_com\_erro + 1$
        \State Salvar log de erro
        \State Continuar para próximo pedido
    \EndIf
\EndFor
\State \Return estatísticas de execução
\end{algorithmic}
\end{algorithm}

\subsection{Características}

\begin{itemize}
    \item \textbf{Ordem de execução}: Fixa, baseada na ordem de chegada (1, 2, 3...)
    \item \textbf{Conflitos de equipamentos}: Não são detectados previamente
    \item \textbf{Gaps temporais}: Não são verificados antecipadamente
    \item \textbf{Função objetivo}: Nenhuma (execução simples)
    \item \textbf{Método de solução}: Greedy (guloso)
    \item \textbf{Complexidade}: $O(n)$ onde $n$ é o número de pedidos
\end{itemize}

\subsection{Componentes do Algoritmo Sequencial}

\begin{table}[h]
\centering
\begin{tabular}{ll}
\toprule
\textbf{Módulo} & \textbf{Função} \\
\midrule
\texttt{executor\_pedidos.py} & Orquestrador da execução sequencial \\
\texttt{gestor\_comandas.py} & Geração de comandas e reservas \\
\texttt{atividades\_modulares.py} & Criação de atividades por pedido \\
\texttt{gestor\_equipamentos.py} & Alocação de equipamentos \\
\bottomrule
\end{tabular}
\caption{Módulos do sistema Sequencial}
\end{table}

\subsection{Vantagens e Desvantagens}

\textbf{Vantagens:}
\begin{itemize}
    \item Implementação simples e direta
    \item Baixo custo computacional
    \item Fácil de depurar e manter
    \item Comportamento previsível
\end{itemize}

\textbf{Desvantagens:}
\begin{itemize}
    \item Não otimiza ordem de execução
    \item Pode violar deadlines desnecessariamente
    \item Ignora conflitos de equipamentos até a execução
    \item Não maximiza utilização de recursos
\end{itemize}

\section{Otimizador PL}

O Otimizador PL é um sistema de otimização de produção que utiliza \textbf{Programação Linear Inteira Mista (MILP)} para determinar a ordem ótima de execução de pedidos.

\subsection{Arquitetura do Otimizador PL}

O sistema é composto por 5 fases principais:

\begin{center}
\begin{tikzpicture}[node distance=1.2cm, auto,
    fase/.style={draw, rectangle, rounded corners, minimum width=5cm, minimum height=0.9cm, font=\small}]

    \node[fase, fill=purple!20] (f0) {Fase 0: Criar Atividades Modulares};
    \node[fase, fill=blue!20, below=of f0] (f1) {Fase 1: Detecção de Modo};
    \node[fase, fill=green!20, below=of f1] (f2) {Fase 2: Cálculo de Horários/Janelas};
    \node[fase, fill=orange!20, below=of f2] (f3) {Fase 3: Otimização PL};
    \node[fase, fill=red!20, below=of f3] (f4) {Fase 4: Execução};

    \draw[->, thick] (f0) -- (f1);
    \draw[->, thick] (f1) -- (f2);
    \draw[->, thick] (f2) -- (f3);
    \draw[->, thick] (f3) -- (f4);
\end{tikzpicture}
\end{center}

\subsection{Componentes do Otimizador PL}

\begin{table}[h]
\centering
\begin{tabular}{ll}
\toprule
\textbf{Módulo} & \textbf{Função} \\
\midrule
\texttt{executor.py} & Orquestrador principal \\
\texttt{detector\_modo.py} & Detecta modo de otimização \\
\texttt{calculador\_horarios\_deterministicos.py} & Cálculo backward scheduling \\
\texttt{calculador\_janelas.py} & Calcula janelas flexíveis \\
\texttt{modelo\_pl\_ordenacao.py} & Modelo PL com OR-Tools CP-SAT \\
\texttt{aplicador\_ordenacao.py} & Executa pedidos (modo determinístico) \\
\texttt{aplicador\_flexivel.py} & Executa pedidos (modo flexível) \\
\bottomrule
\end{tabular}
\caption{Módulos do Otimizador PL}
\end{table}

\subsection{Conceito de $\tau_{max}$ (tempo\_maximo\_de\_espera)}

O parâmetro $\tau_{max}$ define o gap máximo permitido entre atividades consecutivas:

\begin{itemize}
    \item $\tau_{max} = 0$: Atividades consecutivas sem gap (modo DETERMINÍSTICO)
    \item $\tau_{max} > 0$: Gap de até $\tau_{max}$ permitido entre atividades (modo FLEXÍVEL)
\end{itemize}

O modo é detectado automaticamente na Fase 1, após a criação das atividades modulares.

\section{Formulação Matemática}

\subsection{Variáveis de Decisão}

O modelo utiliza variáveis binárias para representar a seleção de pedidos e suas posições:

\begin{equation}
x_{p,j} \in \{0, 1\}
\end{equation}

Onde:
\begin{itemize}
    \item $p$ = índice do pedido ($p \in P$)
    \item $j$ = índice da janela temporal candidata ($j \in J_p$)
    \item $x_{p,j} = 1$ significa que o pedido $p$ será executado na janela $j$
\end{itemize}

Para o modelo de ordenação, também são utilizadas:

\begin{equation}
pos_i \in \{0, 1, \ldots, n-1\}
\end{equation}

Onde $pos_i$ representa a posição do pedido $i$ na ordem de execução.

\subsection{Função Objetivo}

O objetivo é maximizar o número total de pedidos atendidos:

\begin{equation}
\max \sum_{p \in P} \sum_{j \in J_p} x_{p,j}
\end{equation}

\subsection{Restrições}

\subsubsection{Restrição 1: Unicidade por Pedido}

Cada pedido utiliza no máximo uma janela temporal:

\begin{equation}
\sum_{j \in J_p} x_{p,j} \leq 1, \quad \forall p \in P
\end{equation}

\subsubsection{Restrição 2: Permutação Única (Ordenação)}

Para o modelo de ordenação, todas as posições devem ser diferentes:

\begin{equation}
pos_i \neq pos_j, \quad \forall i \neq j
\end{equation}

Esta restrição é implementada usando a constraint \texttt{AddAllDifferent} do OR-Tools.

\subsubsection{Restrição 3: Tempo Máximo de Espera}

Para atividades consecutivas:

\begin{equation}
0 \leq t_{inicio}^{(a+1)} - t_{fim}^{(a)} \leq \tau_{max}
\end{equation}

Quando $\tau_{max} = 0$, isso garante atividades contíguas (modo DETERMINÍSTICO).

\subsubsection{Restrição 4: Conflitos Temporais (Equipamentos)}

Se duas janelas se sobrepõem no tempo, não podem ser selecionadas simultaneamente:

\begin{equation}
x_{p_1,j_1} + x_{p_2,j_2} \leq 1
\end{equation}

Para todo par $(p_1,j_1), (p_2,j_2)$ onde:
\begin{equation}
[t_1^{ini}, t_1^{fim}] \cap [t_2^{ini}, t_2^{fim}] \neq \emptyset
\end{equation}

\subsubsection{Restrição 5: Deadline Rígido}

Para pedidos com deadline obrigatório:

\begin{equation}
x_{p,j} = 0, \quad \forall j : |t_j^{fim} - d_p| > \epsilon
\end{equation}

Onde $d_p$ é o deadline do pedido $p$ e $\epsilon = 1$ minuto (tolerância).

\section{Backward Scheduling}

\subsection{Modo Determinístico ($\tau_{max} = 0$)}

Para o modo determinístico, o sistema utiliza \textbf{backward scheduling} para calcular horários fixos:

\begin{algorithm}
\caption{Backward Scheduling (Determinístico)}
\begin{algorithmic}[1]
\State $t_{atual} \gets deadline$
\For{cada atividade $a$ em ordem reversa}
    \State $t_{fim}^{(a)} \gets t_{atual}$
    \State $t_{inicio}^{(a)} \gets t_{fim}^{(a)} - duracao_a$
    \State $t_{atual} \gets t_{inicio}^{(a)}$
\EndFor
\end{algorithmic}
\end{algorithm}

\subsection{Modo Flexível ($\tau_{max} > 0$)}

Para o modo flexível, o sistema calcula janelas temporais:

\begin{algorithm}
\caption{Cálculo de Janelas Flexíveis}
\begin{algorithmic}[1]
\State $t_{fim}^{tarde} \gets deadline$
\For{cada atividade $a$ em ordem reversa}
    \If{$a$ é última atividade}
        \State $t_{fim}^{cedo}_a \gets deadline$
        \State $t_{fim}^{tarde}_a \gets deadline$
    \Else
        \State $t_{fim}^{tarde}_a \gets t_{fim}^{tarde}$
        \State $t_{fim}^{cedo}_a \gets t_{fim}^{tarde} - \tau_{max}^{sucessora}$
    \EndIf
    \State $t_{inicio}^{tarde}_a \gets t_{fim}^{tarde}_a - duracao_a$
    \State $t_{inicio}^{cedo}_a \gets t_{fim}^{cedo}_a - duracao_a$
    \State $t_{fim}^{tarde} \gets t_{inicio}^{cedo}_a$
\EndFor
\end{algorithmic}
\end{algorithm}

\subsection{Representação Visual}

\textbf{Modo DETERMINÍSTICO ($\tau_{max} = 0$):}
\begin{center}
\begin{tikzpicture}[node distance=0cm]
    \node[draw, rectangle, fill=blue!30, minimum width=2cm, minimum height=0.6cm] (a1) {A1};
    \node[draw, rectangle, fill=blue!30, minimum width=2cm, minimum height=0.6cm, right=of a1] (a2) {A2};
    \node[draw, rectangle, fill=blue!30, minimum width=2cm, minimum height=0.6cm, right=of a2] (a3) {A3};

    \node[below=0.3cm of a1] {\small gap=0};
    \node[below=0.3cm of a2] {\small gap=0};
\end{tikzpicture}
\end{center}

\textbf{Modo FLEXÍVEL ($\tau_{max} = 30min$):}
\begin{center}
\begin{tikzpicture}[node distance=0.5cm]
    \node[draw, rectangle, fill=green!30, minimum width=2cm, minimum height=0.6cm] (a1) {A1};
    \node[draw, rectangle, fill=green!30, minimum width=2cm, minimum height=0.6cm, right=of a1] (a2) {A2};
    \node[draw, rectangle, fill=green!30, minimum width=2cm, minimum height=0.6cm, right=of a2] (a3) {A3};

    \node[below=0.3cm of a1] {\small gap$\leq$30min};
    \node[below=0.3cm of a2] {\small gap$\leq$30min};
\end{tikzpicture}
\end{center}

\subsection{Estratégia de Retry (Modo Flexível)}

Quando uma alocação falha no horário preferencial, o sistema tenta horários anteriores em steps de 5 minutos:

\begin{center}
\begin{tikzpicture}[node distance=0.3cm]
    \node[draw, rectangle, fill=green!30, minimum width=1.5cm, minimum height=0.6cm] (t1) {$t^{tarde}$};
    \node[draw, rectangle, fill=yellow!30, minimum width=1.5cm, minimum height=0.6cm, left=of t1] (t2) {-5min};
    \node[draw, rectangle, fill=yellow!30, minimum width=1.5cm, minimum height=0.6cm, left=of t2] (t3) {-10min};
    \node[left=of t3] (dots) {...};
    \node[draw, rectangle, fill=red!30, minimum width=1.5cm, minimum height=0.6cm, left=of dots] (tn) {$t^{cedo}$};

    \draw[->, thick] (t1) -- (t2);
    \draw[->, thick] (t2) -- (t3);
    \draw[->, thick] (t3) -- (dots);
    \draw[->, thick] (dots) -- (tn);

    \node[above=0.3cm of t1] {\small Preferencial};
    \node[above=0.3cm of tn] {\small Limite};
\end{tikzpicture}
\end{center}

\section{Complexidade Computacional}

\begin{table}[h]
\centering
\begin{tabular}{ll}
\toprule
\textbf{Componente} & \textbf{Complexidade} \\
\midrule
Variáveis de decisão & $O(n \cdot k)$ \\
Restrições de unicidade & $O(n)$ \\
Restrições de conflito & $O(n^2 \cdot k^2)$ (pior caso) \\
\bottomrule
\end{tabular}
\caption{Complexidade por componente ($n$ = pedidos, $k$ = janelas/pedido)}
\end{table}

O solver OR-Tools CP-SAT utiliza técnicas avançadas de propagação de restrições e busca para encontrar soluções ótimas ou viáveis dentro do timeout configurado (padrão: 600 segundos).

\section{Comparação: Sequencial vs Otimizado}

\begin{table}[h]
\centering
\begin{tabular}{lcc}
\toprule
\textbf{Aspecto} & \textbf{Sequencial} & \textbf{Otimizado PL} \\
\midrule
Ordem de execução & Fixa & Otimizada \\
Conflitos de equipamentos & Ignora & Modela \\
Modo DETERMINÍSTICO ($\tau_{max}=0$) & - & Sim \\
Modo FLEXÍVEL ($\tau_{max}>0$) & - & Sim \\
Detecção automática de modo & Não & Sim \\
Janelas flexíveis & Não & Sim (modo FLEXÍVEL) \\
Retry na alocação & Não & Sim (modo FLEXÍVEL) \\
Função objetivo & Nenhuma & Max pedidos \\
Método de solução & Greedy & MILP \\
Complexidade & $O(n)$ & $O(n^2 \cdot k^2)$ \\
\bottomrule
\end{tabular}
\caption{Comparação entre métodos}
\end{table}

\section{Implementação}

\subsection{Exemplo de Uso - Sequencial}

\begin{lstlisting}[language=Python, basicstyle=\ttfamily\small]
from services.gestores.producao import ExecutorPedidos

# Criar executor
executor = ExecutorPedidos()

# Executar pedidos sequencialmente
sucesso = executor.executar_sequencial(pedidos_convertidos)

# Verificar estatisticas
stats = executor.obter_estatisticas()
print(f"Executados: {stats['pedidos_executados']}")
print(f"Com erro: {stats['pedidos_com_erro']}")
\end{lstlisting}

\subsection{Exemplo de Uso - Otimizador PL}

\begin{lstlisting}[language=Python, basicstyle=\ttfamily\small]
from otimizador import Executor

# Criar executor
executor = Executor()
executor.inicializar()

# Otimizar pedidos (detecta modo automaticamente)
solucao = executor.otimizar_pedidos(pedidos)

# Verificar resultados
print(f"Modo detectado: {solucao.modo_detectado}")
print(f"Pedidos atendidos: {solucao.pedidos_atendidos}")
print(f"Taxa de sucesso: {solucao.estatisticas['taxa_sucesso']}%")
\end{lstlisting}

\subsection{Parâmetros de Configuração}

\begin{itemize}
    \item \textbf{timeout\_segundos}: Tempo máximo para resolução (padrão: 600s)
    \item \textbf{inicio\_jornada}: Horário de início da jornada de trabalho
\end{itemize}

\section{Quando Usar Cada Algoritmo}

\subsection{Use o Algoritmo Sequencial quando:}

\begin{itemize}
    \item Os pedidos não possuem deadlines rígidos
    \item A ordem de chegada é importante para o cliente
    \item O número de pedidos é pequeno (< 10)
    \item Simplicidade é mais importante que otimização
    \item Recursos computacionais são limitados
\end{itemize}

\subsection{Use o Otimizador PL quando:}

\begin{itemize}
    \item Existem deadlines rígidos a serem respeitados
    \item Há conflitos potenciais de equipamentos
    \item É necessário maximizar o número de pedidos atendidos
    \item O número de pedidos é grande (> 10)
    \item As atividades podem ter $\tau_{max} > 0$ (gaps permitidos)
    \item É necessário flexibilidade na alocação temporal
    \item O sistema deve adaptar-se a conflitos de recursos
\end{itemize}

\section{Conclusão}

O sistema SIVIRA oferece dois algoritmos para alocação de pedidos:

\begin{enumerate}
    \item O \textbf{Algoritmo Sequencial} é ideal para cenários simples, onde a ordem de chegada deve ser respeitada e não há restrições temporais complexas. Sua simplicidade garante baixo custo computacional e comportamento previsível.

    \item O \textbf{Otimizador PL} oferece uma abordagem matematicamente fundamentada para o problema de scheduling de produção, suportando ambos os modos:
    \begin{itemize}
        \item \textbf{DETERMINÍSTICO} ($\tau_{max} = 0$): Atividades contíguas, sem gaps
        \item \textbf{FLEXÍVEL} ($\tau_{max} > 0$): Janelas temporais com retry automático
    \end{itemize}
    O modo é detectado automaticamente antes da execução, garantindo comportamento correto em ambos os casos.
\end{enumerate}

A arquitetura modular permite fácil extensão para novos modos de operação e integração com o sistema de menu existente.

\end{document}
